% ----------------------
\subsection{The structural fallacy}
% ----------------------
	% the fallacy of inferring that, since God appears to have designed something a certain way, then therefore that way reflects a true aspect of the plan of salvation, or some aspect of eternity.

	% -----
	\subsubsection{Teleological explanations and the teleological fallacy}
	% -----
	
		In the previous essay,
			I discussed teleological explanations and the teleological fallacy.
		A teleological explanation is any explanation that makes sense of an object or event by referring to its purpose,
			design,
				or the reason for which it was created.
		So say I'm driving along one day and see someone in front of me with their right blinker on.
		I explain what they're doing teleologically by referring to the purpose they must have in mind:
			they're going to turn right.
			
		Teleological explanations work in domains where there are designs and reasons.
		Human behavior is one such area;
			biology is another.
		But teleological explanations fail or at least aren't usually appropriate in domains where there aren't designs and reasons.
		In the last essay I mentioned that astronomy is an example of such a domain.
		
		That point about teleological explanations led us to the \textit{teleological fallacy}:
			the mistake of inferring that,
				because something \textit{is} a certain way,
					it must have been \textit{designed} to be that way.
		We're particularly susceptible to this fallacy when reasoning about the gospel,
		and so I also distinguished the \textit{gospel teleological fallacy}:
			the mistake of inferring that,
				because something \textit{is} a certain way in the gospel,
					God must have \textit{designed} it to be that way.
		I gave some examples of things we might reason about where this fallacy would be tempting.
		Like I said before,
			I'll be covering some of those things in greater detail later on.
			
		In this essay,
			though,
				I'm going to extend the gospel teleological fallacy into (yet) another one,
					which I'll call the \textit{structural fallacy},
						or \textit{gospel structural fallacy}.
		I'll state it here and then discuss it:
		
		\begin{description}
				\item[The Gospel Structural Fallacy] The gospel structural fallacy is the fallacy of inferring that, since God appears to have designed something in a certain manner, that manner therefore reflects a true aspect of the plan of salvation's structure, or some aspect of the structure of eternity.
			\end{description}
	
	% -----
	\subsubsection{Explaining the structural fallacy}
	% -----
	
		The second half of the structural fallacy may be a bit tricky to understand.
		By ``a true aspect of the plan of salvation's structure,''
			I mean some fact about the way the plan of salvation is set up:
				for example,
					that there are three kingdoms of glory,
						or that a person has to be sealed in order to enter into the highest degree of the celestial kingdom.
		I think of facts like those as facts about the \textit{structure} of the plan of salvation,
			or again,
				about the way it is set up or proceeds.
		By ``some aspect of the structure of eternity'' I mean the same thing,
			except that this latter way of putting it pushes us beyond the bounds of the plan of salvation to think of how things are set up in the very very long run,
				such as after a person is exalted.
		If for some reason an exalted person must eat cake for breakfast every day,
			then the mandatory cake-eating would be part of the \textit{structure} of eternity.
		More serious parts of eternity's structure would (probably) include things like family relationships and having power to control matter in certain ways.
		
		I do still hope there's cake though.
	
	% -----
	\subsubsection{The teleological fallacy and the structural fallacy}
	% -----
	
		Now let's think about how the structural fallacy works.
		First,
			look at it in relation to the teleological fallacy.
		The gospel teleological fallacy makes one inference:
			from the fact that something is a certain way,
				to the claim that God must have designed it to be that way.
		The structural fallacy then picks up where the teleological fallacy leaves off---%
			it begins from the claim that God must have designed the thing to be that way.
		Then it makes one more inferential step and concludes that,
			because God designed the thing that way,
				the thing therefore reflects something about the plan of salvation,
					the structure of eternity,
						or maybe even the universe itself.
		
		In the previous essay I gave a short list of things we might be tempted to commit the teleological fallacy about.
		One was the fact that men bear the priesthood,
			but not women.
		So someone might make the following inference:
		
		\begin{table}[H]
		\centering
		\begin{tabular}{p{7cm}p{0.5cm}p{7cm}}
		Men bear the priesthood, but not women. & $\rightarrow$ & Therefore, God intended (designed, purposed) only men to bear the priesthood, and not women.
		\end{tabular}
		% \caption{Caption}
		\end{table}
		
		(The $\rightarrow$ symbol means something like ``Therefore'' or ``It follows that.'')
		That first diagram just illustrates the teleological fallacy.
		The structural fallacy then goes further,
			and in diagram form it would look like this:
		
		\begin{table}[H]
		\centering
		\begin{tabular}{p{7cm}p{0.5cm}p{7cm}}
		God intended only men to bear the priesthood, and not women. & $\rightarrow$ & Therefore, it is a feature of the plan of salvation in the long run that only men bear the priesthood, and not women (or, it will always be true that only men bear the priesthood, and not women).
		\end{tabular}
		% \caption{Caption}
		\end{table}
		
		When we compare these two fallacies,
			we can see that the structural fallacy begins where the teleological fallacy ends.
		It then takes one more step to conclude something about the nature of eternity itself.
		The rules of logic%
			\footnote{This rule is called ``hypothetical syllogism.''}
			allow us to put the two fallacies together into one,
				like this:
				
		\begin{table}[H]
		\centering
		\begin{tabular}{p{7cm}p{0.5cm}p{7cm}}
		Men bear the priesthood, and not women. & $\rightarrow$ & Therefore, it is a feature of the plan of salvation in the long run that only men bear the priesthood, and not women (or, it will always be true that only mean bear the priesthood, and not women).
		\end{tabular}
		% \caption{Caption}
		\end{table}
		
		This sort of reasoning goes from an observation of the way things are,
			to a conclusion about what they are like in the very long run,
				or what they are like in eternity.
		In this representation we've just skipped the middle step of teleological reasoning.
				
	% -----
	\subsubsection{What justifies a structural inference?}
	% -----
	
		So far I've only been referring to the structural \textit{fallacy},
			but we should remember that not every structural inference has to be a fallacy.
		Here's one that may not be a fallacy:
		
		\begin{table}[H]
		\centering
		\begin{tabular}{p{7cm}p{0.5cm}p{7cm}}
		We live in family units on Earth. & $\rightarrow$ & Therefore, it is a feature of the plan of salvation in the long run that we will live in family units (or, it will always be true that we live in family units).
		\end{tabular}
		% \caption{Caption}
		\end{table}
		
		Suppose this one isn't a fallacy (I'm not sure whether it is).
		If it isn't a fallacy,
			then it's just an example of what we could call structural \textit{reasoning} or a structural \textit{inference}.
		Just as there are instances of teleological reasoning which succeed or are good,
			there are probably instances of structural reasoning which are good too.
		Not every example of structural reasoning needs to be a fallacy.
		
		In the previous essay on teleological reasoning,
			we discussed how teleological inferences might fail.
		In particular,
			they're risky and prone to failure in domains where no design processes have been acting,
				like astronomy.
		No design processes means no intention-based reasons,
			and so teleological reasoning will probably fail.
		In domains that do have reasons,
			teleological inferences might be good ones.
		We could also say that they would be \textit{justified} in such a domain,
			or that the presence of reasons \textit{justifies} the general use of teleological reasoning.
			
		In the same way,
			we can also ask what justifies a structural inference.
		Teleological inferences need reasons,
			but what do structural inferences need?
		Let's remind ourselves of the structural inference we saw earlier,
			which begins with a teleological statement:
			
		\begin{table}[H]
		\centering
		\begin{tabular}{p{7cm}p{0.5cm}p{7cm}}
		God intended only men to bear the priesthood, and not women. & $\rightarrow$ & Therefore, it is a feature of the plan of salvation in the long run that only men bear the priesthood, and not women (or, it will always be true that only men bear the priesthood, and not women).
		\end{tabular}
		% \caption{Caption}
		\end{table}
		
		What would it take for this inference to be a \textit{good} one---%
			what would it take to \textit{justify} it?
		The left side is a statement about God's intention or design,
			and the whole thing asks us to believe the conclusion (the right side) \textit{on the basis of} the left side.
		In other words,
			this inference presents a conclusion and offers the left side as the evidence for that conclusion.
		Because God intended only men to bear the priesthood,
			on that basis we are supposed to have evidence that it will always be true that only men will bear the priesthood.
			
		Here's one way of explaining how this inference could be good,
			or how the statement on the left could really be evidence for the one on the right:
				the inference would be good when God looks at the structure of eternity,
					and then decides his intention based on that structure.
		So imagine God looks down the road into eternity and sees that only men will be able to bear the priesthood.
		Recognizing this aspect of eternity,
			God therefore purposes to only give men the priesthood.
		If this were true,
			the inference from God's intention to the structure of eternity would be a good one.
		The justification would be that God is acting according to the constraints he perceives,
			and so our inferences involving those constraints are reliable.
			
		That isn't the only way that the inference could be good,
			though. 
		It could also be good because,
			rather than being constrained by the structure of eternity,
				God just \textit{wants} only men to have the priesthood.
		If that were true,
			it would be easy to see why the inference is a good one---%
				if we knew that God wanted only men to have the priesthood,
					and that he intended only men to have the priesthood by giving it only to them,
						it would be reasonable to conclude that only men will have it in the long run.
		The structural inference would be justified.
		
		Notice that we've just identified two \textit{different} ways of justifying the \textit{same} inference.
		One way requires an assumption about some constraint that God sees;
			another way requires an assumption about God's will.
		The inference is good on either assumption,
			but it's good for a different reason.
		That point is really important,
			and I'll talk about it more in the next essay.
	
	% -----
	\subsubsection{Some (possible) examples of the structural fallacy}
	% -----
	
		So now we know when a structural inference might be good and not just a fallacy.
		Let's mention some other possible examples.
		
		Take the fact that very few women appear in the LDS canon of scripture,
			and very \textit{very} few in restoration scripture,
				such as the Book of Mormon.
		Someone---%
			not you,
				of course---%
					could reason like this:
						``I observe that few women appear in restoration scripture.
		It therefore must be that God intended for few women to appear.
		On that basis,
			it must therefore be that women don't play a large role in the plan of salvation (or eternity) at large.''
		In that line of reasoning,
			the first ``therefore'' marks a teleological inference,
				and the second ``therefore'' marks a structural inference.
				
		Now that inference won't look very good to most people.
		Fundamentalists might like it,
			but everyone else will recognize right away that there's a simple way to challenge the inference---%
				you just give the obvious alternative explanation.
		If you're a fundamentalist and it has not occurred to you yet,
			it's that the writers of the scriptures were all men,
				and so tended to focus on themselves or other men involved in the events they wrote about.
		Since we know the writers of the scriptures were men,
			if it were also true that they tended to have a biased focus,
				then that line of reasoning wouldn't be very strong.
				
		We'll call that alternative explanation a \textit{naturalistic} one,
			because it refers only to the natural order of things,
				or how the natural world is (people are part of the natural world,
					in case you forgot).
		For most people,
			the naturalistic explanation will be preferable,
				because otherwise you attribute to God some weird inclination to disappear women from the scriptures,
					and why would he do that?
		It wouldn't make much sense,
			so the alternative explanation has a lot of appeal.
		(In keeping with the pattern I've already set,
			let me mention here that I'll have more to say about naturalistic explanations in the future.)
			
		Another possible example is the priesthood organization of the church.
		There is a president,
			a first presidency,
				a quorum of twelve apostles,
					several quorums of seventy,
						presidents of those seventies,
							and so on.
		Then you've got auxiliary presidencies and everything else.
		Observing this organization,
			one might infer that God intended things to organized that way,
				and then might further infer that things will \textit{always} be like that---%
					that there will always be a president,
						apostles, and so on.
		
		For some people,
			this inference will look more attractive than the one about women and the scriptures.
		If you think of heads of families as occupying positions of priesthood leadership (I'm referring to both men and women here),
			then it might seem even more reasonable.
		But it might also be an instance of the structural fallacy.
		
		One place we'll be sure \textit{not} to commit the fallacy is in thinking about the veil of forgetfulness,
			which we've mentioned before.
		Even if it were true that God intended for us to have the veil,
			it wouldn't follow that the veil is a feature of eternity.
		It's supposed to be temporary,
			after all.
		We aren't prone to committing the structural fallacy here,
			because we have other reasons for thinking that the ``structure'' isn't permanent.
	
	% -----
	\subsubsection{The teleological and structural fallacies: locating the error}
	% -----
	
		Two more sections in this essay,
			and then we'll move on to the next.
		
		First,
			a word about the teleological and structural fallacies together.
		So far we've seen many examples of a common pattern of reasoning.
		It goes like this:
		
		\begin{table}[H]
		\centering
		\begin{tabular}{p{4cm}p{0.5cm}p{4cm}p{0.5cm}p{4cm}}
		Observation & $\rightarrow$ & Teleological statement & $\rightarrow$ & Structural statement
		\end{tabular}
		% \caption{Caption}
		\end{table}
		
		The pattern starts with some observation relevant to the gospel.
		It then moves from that observation to a teleological statement,
			or a statement involving God's intention or design or will.
		Then it goes from there to a structural statement,
			about how things are or will be in eternity.
		
		Here's the reason I've been so careful in describing this reasoning:
			\textit{most errors involving the structural fallacy are really errors involving the teleological fallacy}.
		That is,
			if you've got a structural conclusion that seems wrong,
				odds are the bad inference occurs not in going from the teleological statement to the structural one,
					but in going from the observation to the teleological statement.
		When you make that error to start with,
			the structural fallacy is then like making one more mistake on top of another.
		But the original and most serious mistake is \textit{not} in going from the teleological point to the structural one---%
			it happens earlier than that.
			
		On the other hand,
			however,
				most of the time we do happen to \textit{notice} the error in reasoning when we go from the teleological to the structural.
		That's because the further inference tends to draw out the true consequences of the teleological inference,
			and when we find the consequences unacceptable,
				we reject the inference.
		But we often reject it at the step going from the teleological to the structural,
			when the real mistake happens \textit{before} that.
		The real error is the teleological fallacy,
			and \textit{that} is where we need to reject the inference---%
				we shouldn't even allow the inferential chain to get started.
		But if we don't separate them out as I have here,
			we won't see where the real error is.
		We'll just think it's wherever we happened to notice it.
	
		So take this inference again:
		
		\begin{table}[H]
		\centering
		\begin{tabular}{p{7cm}p{0.5cm}p{7cm}}
		Men bear the priesthood, and not women. & $\rightarrow$ & Therefore, it is a feature of the plan of salvation in the long run that only men bear the priesthood, and not women (or, it will always be true that only mean bear the priesthood, and not women).
		\end{tabular}
		% \caption{Caption}
		\end{table}
		
		We start with the observation and end with a structural statement.
		If we didn't like that conclusion,
			we might think that the argument committed the structural fallacy.
		But we now know that the argument as set up here hides its own middle step:
		
		\begin{table}[H]
		\centering
		\begin{tabular}{p{7cm}p{0.5cm}p{7cm}}
		Men bear the priesthood, but not women. & $\rightarrow$ & Therefore, God intended (designed, purposed) only men to bear the priesthood, and not women.
		\end{tabular}
		% \caption{Caption}
		\end{table}
		
		The middle step is the teleological one.
		I'm suggesting that the true error here is at the teleological step,
			not at the structural one.
		In other words,
			we shouldn't even make the teleological inference.
		But we can't see that unless we've distinguished between the two steps,
			and between their corresponding fallacies,
				well enough to tell them apart and understand how they work together.
		
	% -----
	\subsubsection{Why the teleological and structural fallacies are important}
	% -----
	
		The next and final point is on why these fallacies are so important to recognize.
		
		The answer is simple:
			in the LDS church,
				both of them are \textit{extremely} common!
		There is so much organization,
			so much bureaucracy,
				so many programs,
					and so much of everything else,
						almost all backed by leaders considered to be inspired,
							that you start to think there must really be something to all that organization after all.
		You start to think the programs or the hierarchy are modeled after some higher or more lasting reality,
			such as an aspect of eternity or some deep part of the plan of salvation.
		I have even heard people speculate,
			in all seriousness,
				about whether there would be home teaching in the celestial kingdom!
				
		Religious believers in general can sometimes be guilty of finding meaning and purpose where there is none,
			and given their belief in God,
				it's easy for them to attribute those meanings and purposes to God.
		This gives the purposes they've discerned a more lasting feel,
			as though they were latching onto deeper principles.
		And LDS believers,
			with their expanded canon of scripture,
				array of local and international leaders,
					and emphasis on personal revelation,
						may find it even harder to avoid these fallacies.
		
		The first step to not making these errors is to identify them and bring them out in the open.
		That's what I've done here and in the previous essay.
		Seeing the fallacies in action also highlights the need for alternative modes of explanation in understanding gospel phenomena.
		I'll begin exploring one of those alternative modes next.